\chapter{Le sette specie di settima}
\label{chap:sette_specie}

Dopo aver studiato la settima di dominante e i suoi derivati diretti, conviene considerare l'intera famiglia delle quadriadi. In questo capitolo sistematizzeremo le diverse "specie" di settima, ovvero le diverse combinazioni di terze sovrapposte che possono nascere sui gradi della scala.

Per quanto riguarda il comportamento generale delle dissonanze, la loro preparazione e la loro risoluzione obbligatoria, si rimanda integralmente al \textbf{Capitolo \ref{chap:teoria_dissonanza}}.

\section{Classificazione generale}
\label{sec:classificazione_specie}

Tradizionalmente si distinguono sette specie di accordi di settima, determinate dalla natura della triade di base e dalla qualità dell'intervallo di settima rispetto alla fondamentale.

\begin{enumerate}
    \item \textbf{Prima specie:} Triade Maggiore + settima minore (Settima di dominante).
    \item \textbf{Seconda specie:} Triade minore + settima minore.
    \item \textbf{Terza specie:} Triade diminuita + settima minore (Settima semidiminuita).
    \item \textbf{Quarta specie:} Triade Maggiore + settima Maggiore.
    \item \textbf{Quinta specie:} Triade diminuita + settima diminuita (Settima diminuita).
    \item \textbf{Sesta specie:} Triade minore + settima Maggiore.
    \item \textbf{Settima specie:} Triade aumentata + settima Maggiore.
\end{enumerate}

\begin{esempio}[Le sette specie costruite su Do]
	\label{ex:sette_specie}
	Le quattro prime specie e le tre restanti, costruite sulla stessa fondamentale. Sotto ogni accordo sono indicate le qualità delle tre terze sovrapposte (M = maggiore, m = minore).
	\musica{es_cap13_10}
	\musica{es_cap13_11}
\end{esempio}

In un corso di armonia complementare ci concentreremo sulle casistiche più comuni per l'analisi e la composizione standard, tralasciando le specie più rare, che negli esami non trovano applicazione.

\section{Collocazione delle specie sui gradi delle scale}
\label{sec:collocazione_scale}

Gli accordi di settima si formano sovrapponendo terze \textbf{restando nei suoni della scala}: a ogni grado corrisponde quindi un accordo di specie determinata. Le sigle sotto gli accordi indicano il grado (cifre romane) e la specie.

\begin{esempio}[Settime sui gradi di Do maggiore]
	\label{ex:scala_do}
	\musica{es_cap13_12}
\end{esempio}

\begin{esempio}[Settime sui gradi di La minore naturale]
	\label{ex:scala_lan}
	Nel minore naturale la dominante è \emph{minore} (\sigla{V}{7}{} di 2\textsuperscript{a} specie) e manca la sensibile: per questo nella pratica tonale si ricorre alla scala armonica.
	\musica{es_cap13_13}
\end{esempio}

\begin{esempio}[Settime sui gradi di La minore armonica]
	\label{ex:scala_la_arm}
	Sono gli accordi caratteristici del modo minore: \sigla{V}{7}{} (dominante, con la sensibile Sol\diesis), \sigla{VII}{7}{} (settima diminuita di sensibile) e \sigla{II}{7}{} (semidiminuita, preparazione alla dominante). Gli accordi con quinta eccedente (\sigla{III}{7}{}) e minore-maggiore (\sigla{I}{7}{}) sono rari.
	\musica{es_cap13_14}
\end{esempio}

\begin{esempio}[Settime sui gradi di La minore melodica ascendente]
	\label{ex:scala_la_mel}
	\musica{es_cap13_15}
\end{esempio}

\begin{center}
\begin{tabular}{|c|c|c|c|c|}
\hline
Grado & Do maggiore & La min.\ naturale & La min.\ armonica & La min.\ melodica asc. \\
\hline
I & 4\textsuperscript{a} & 2\textsuperscript{a} & 6\textsuperscript{a} & 6\textsuperscript{a} \\
II & 2\textsuperscript{a} & 3\textsuperscript{a} & 3\textsuperscript{a} & 2\textsuperscript{a} \\
III & 2\textsuperscript{a} & 4\textsuperscript{a} & 7\textsuperscript{a} & 7\textsuperscript{a} \\
IV & 4\textsuperscript{a} & 2\textsuperscript{a} & 2\textsuperscript{a} & 1\textsuperscript{a} \\
V & 1\textsuperscript{a} & 2\textsuperscript{a} & 1\textsuperscript{a} & 1\textsuperscript{a} \\
VI & 2\textsuperscript{a} & 4\textsuperscript{a} & 4\textsuperscript{a} & 3\textsuperscript{a} \\
VII & 3\textsuperscript{a} & 1\textsuperscript{a} & 5\textsuperscript{a} & 3\textsuperscript{a} \\
\hline
\end{tabular}
\end{center}

\section{Prima e Terza specie: richiami}
\label{sec:richiami_specie}

Le specie più importanti sono già state trattate nei capitoli precedenti:

\begin{itemize}
    \item \textbf{Prima Specie (Settima di dominante):} Si rimanda al \textbf{Capitolo \ref{chap:settima_dominante}}. Non richiede la preparazione obbligatoria della settima (come la terza e la quinta specie).
    \item \textbf{Terza Specie (Settima di sensibile / semidiminuita):} Si rimanda al \textbf{Capitolo \ref{chap:vii7_ii7}}. Ricordiamo che, sulla settima di sensibile del VII grado del modo maggiore, la preparazione della settima \textbf{non è obbligatoria}; lo è invece per l'accordo sul \textbf{II grado del modo minore} (fase di consonanza nell'accordo precedente). In ogni caso la settima risolve scendendo di grado. La sua funzione è di dominante (sul VII grado del maggiore) o di sottodominante (sul II grado del minore).
\end{itemize}

\section{Seconda Specie}
\label{sec:seconda_specie_dettaglio}

La seconda specie è formata da una \textbf{triade minore} e una \textbf{settima minore}. L'accordo di gran lunga più importante di questa specie è quello costruito sul \textbf{secondo grado (II)} del modo maggiore.

\subsection{Funzione e risoluzione}
La sua funzione è tipicamente di \textbf{sottodominante}. Essendo una dissonanza di seconda specie, la settima deve essere rigorosamente \textbf{preparata}. Risolve naturalmente sulla dominante (V o \sigla{V}{7}{}).

\begin{regola}[Collegamento I-\sigla{II}{7}{}-V e raddoppio della terza]
    \label{reg:ii7_raddoppio}
    Nel collegamento tra il I grado e il \sigla{II}{7}{} (e la successiva risoluzione su V), è frequente e consigliabile \textbf{raddoppiare la terza} dell'accordo di \sigla{II}{7}{} (che corrisponde al IV grado della scala). Serve a evitare \textbf{quinte parallele} tra la fondamentale e la quinta del II grado nel movimento verso l'accordo di dominante.

	In alternativa (o in aggiunta) si può \textbf{omettere la quinta} del \sigla{II}{7}{}, raddoppiando una delle altre voci (es. la fondamentale) e lasciando fondamentalmente la terza raddoppiata, come nell'esempio seguente.
\end{regola}

\begin{esempio}[\sigla{II}{7}{} in Do Maggiore]
    \musica{es_cap13_minore_do_1}
    \musica{es_cap13_minore_do_2}
    Il \sigla{II}{7}{} risolve sul V: la settima Do scende a Si (sensibile), che non si raddoppia. Il primo rivolto (\sigla{II}{6}{5}, basso Fa) sale a Sol, il secondo (\sigla{II}{4}{3}, basso La) scende a Sol, il terzo (\sigla{II}{2}{}, basso Do) scende a Si sul \sigla{V}{6}{}.
\end{esempio}

\begin{esempio}[Seconda specie in La minore: \sigla{I}{7}{} $\rightarrow$ IV]
	\label{ex:i7_minore}
	Nel minore naturale la settima minore sul I grado (La-Do-Mi-Sol) risolve sul IV (Re-Fa-La): la settima Sol scende a Fa, terza del IV. Il comportamento è identico a quello dell'esempio in Do maggiore.
	\musica{es_cap13_minore_la_1}
	\musica{es_cap13_minore_la_2}
\end{esempio}

\subsection{Collocazione sulle scale}
\label{subsec:ii7_collocazione}

\begin{regola}[Gradi della seconda specie]
	\label{reg:ii7_gradi}
	L'accordo di settima minore (seconda specie) si trova:
	\begin{itemize}
		\item Sul \textbf{II}, \textbf{III} e \textbf{VI} grado della scala maggiore;
		\item Sul \textbf{IV} grado della scala minore armonica.
	\end{itemize}
	L'accordo di gran lunga più importante è quello sul \textbf{II grado del modo maggiore}, con funzione di sottodominante.
\end{regola}

\subsection{Stato fondamentale}
\label{subsec:ii7_stato_fondamentale}

\begin{regola}[\sigla{II}{7}{} allo stato fondamentale]
	\label{reg:ii7_stato_fondamentale}
	Al II grado al basso, la settima va \textbf{preparata} nell'accordo precedente (presente nella stessa voce, con durata pari o superiore) e \textbf{tenuta}, poi risolve \textbf{scendendo per grado congiunto}. Le altre note, compresa la fondamentale, sono \textbf{parti libere} e si possono raddoppiare. La quinta si può omettere raddoppiando una delle altre voci (mai la settima).
\end{regola}

Risolve naturalmente sul \textbf{V grado}, armonizzato con la triade semplice, con la cadenza composta consonante o dissonante, oppure con la $\sigla{V}{7}{}$.

\subsection{I tre rivolti}
\label{subsec:ii7_rivolti}

\begin{center}
\begin{tabular}{|l|l|l|l|l|}
\hline
Stato & Cifratura & Grado al basso & Risolve su & Basso obbligato \\
\hline
Fond. & 7 & II & V (o \sigla{V}{7}{}) & No \\
I rv. & 6/5 & IV & V & Sì (sale) \\
II rv. & 4/3 & VI & V & Sì (scende) \\
III rv. & 4/2 & I (settima) & \sigla{V}{6}{} o \sigla{V}{6}{5} & Sì (scende) \\
\hline
\end{tabular}
\end{center}

\begin{itemize}
	\item \textbf{\sigla{II}{6}{5}} (IV grado al basso): il basso risolve per grado congiunto ascendente sulla dominante. È impiegato diffusamente nella formula cadenzale $I$ (oppure $\sigla{I}{6}{}$ / $VI$) $- \sigla{II}{6}{5} - \sigla{I}{6}{4} - V$. Può risolvere anche sulla triade o settima di sensibile allo stato di rivolto.
	\item \textbf{\sigla{II}{4}{3}} (VI grado al basso): il basso risolve per moto discendente sul V grado.
	\item \textbf{\sigla{II}{2}{}} (I grado al basso, la settima): la settima al basso deve essere preparata e risolve per moto discendente sul VII grado, armonizzato con il primo rivolto della dominante ($\sigla{V}{6}{}$ o $\sigla{V}{6}{5}$).
\end{itemize}

\begin{osservazione}[Progressioni e IV grado del modo minore]
	\label{oss:ii7_progressioni}
	\begin{itemize}
		\item Il \sigla{VI}{7}{} e gli altri gradi di seconda specie si usano scolasticamente nelle \textbf{progressioni discendenti} (es. $\sigla{III}{7}{} - \sigla{VI}{7}{} - \sigla{II}{7}{}$).
		\item La seconda specie sul IV grado del modo minore è un accordo \textbf{derivato} (da una nona di terza specie sul II grado senza fondamentale) e si comporta esattamente come la quarta specie sul IV grado del modo maggiore.
	\end{itemize}
\end{osservazione}

\section{Quarta Specie}
\label{sec:quarta_specie_dettaglio}

La quarta specie è formata da una \textbf{triade Maggiore} e una \textbf{settima Maggiore}. Nel corso di armonia studieremo solo l'accordo costruito sul \textbf{primo grado (I)} del modo maggiore.

\subsection{Uso e risoluzione}
Questo accordo ha una tensione particolare e risolve generalmente sulla \textbf{sottodominante (IV)}. La settima (che è la sensibile della scala) deve essere preparata e risolve scendendo di grado verso la terza del IV grado.

\begin{esempio}[\sigla{I}{7}{} della quarta specie]
    \musica{es_cap13_maggiore_do_1}
    \musica{es_cap13_maggiore_do_2}
    Il \sigla{I}{7}{} risolve sul IV: la settima Si scende a La. Primo rivolto (basso Mi) sale a Fa; secondo rivolto (basso Sol) scende a Fa; terzo rivolto (basso Si, la settima) scende a La sul \sigla{IV}{6}{}.
\end{esempio}

\begin{esempio}[Quarta specie in La minore: \sigla{VI}{7}{} $\rightarrow$ II]
    \label{ex:vi7_minore}
    Sul VI grado della minore armonica la settima maggiore (Fa-La-Do-Mi) risolve per quinta discendente sul II grado (triade diminuita Si-Re-Fa): la settima Mi scende a Re.
    \musica{es_cap13_maggiore_la_1}
    \musica{es_cap13_maggiore_la_2}
\end{esempio}

\subsection{Collocazione sulle scale}
\label{subsec:iv7_collocazione}

\begin{regola}[Gradi della quarta specie]
	\label{reg:iv7_gradi}
	L'accordo di settima maggiore (quarta specie) si trova:
	\begin{itemize}
		\item Sul \textbf{I} e \textbf{IV} grado della scala maggiore;
		\item Sul \textbf{VI} grado della scala minore armonica;
		\item Sul \textbf{III} e \textbf{VI} grado della scala minore melodica discendente.
	\end{itemize}
	L'analisi del corso si concentra sul \textbf{I grado del modo maggiore}.
\end{regola}

\subsection{Rivolti}
\label{subsec:iv7_rivolti}

Al I grado, l'accordo di quarta specie ha funzione di \textbf{sottodominante} (o di settima sul IV grado). La settima va preparata nella voce; le altre note sono parti libere e la quinta si può omettere a favore del raddoppio della fondamentale.

\begin{center}
\begin{tabular}{|l|l|l|l|l|}
\hline
Stato & Cifratura & Grado al basso & Risolve su & Basso obbligato \\
\hline
Fond. & 7 & I & IV (o \sigla{IV}{7}{}) & No \\
I rv. & 6/5 & III & IV & Sì (sale) \\
II rv. & 4/3 & V & IV & Sì (scende) \\
III rv. & 4/2 & VII (settima) & \sigla{IV}{6}{} & Sì (scende) \\
\hline
\end{tabular}
\end{center}

\section{Quinta Specie}
\label{sec:quinta_specie_dettaglio}

La quinta specie è la \textbf{settima diminuita} (triade diminuita + settima diminuita). È l'accordo più dissonante e simmetrico.

\subsection{Natura e collocazione}
\label{subsec:vii_dim7_natura}

\begin{regola}[Grado della quinta specie]
	\label{reg:vii_dim7_gradi}
	La settima diminuita (quinta specie) si trova sul \textbf{VII grado della scala minore armonica} (nella melodica ascendente il VII grado dà invece una semidiminuita), dove è chiamata \textbf{settima di sensibile del modo minore}. È l'unico accordo di settima a contenere l'intervallo di \textbf{settima diminuita}, perciò è chiaramente distinguibile all'ascolto.
\end{regola}

È un accordo \textbf{principale e derivato}: nasce dall'accordo di nona di dominante minore senza la fondamentale, e per questo \textbf{non necessita della preparazione} della settima.

\subsection{Uso sul VII grado del modo minore}
Si studia principalmente sul \textbf{VII grado del modo minore} (sulla sensibile), dove svolge funzione di dominante estrema. Alla risoluzione sul I grado valgono le norme della settima di sensibile di terza specie (Capitolo \ref{chap:vii7_ii7}): la sensibile sale alla tonica, la quinta e la settima scendono per grado.

\begin{esempio}[\sigla{VII°}{7}{} $\rightarrow$ I in La Minore]
	\label{ex:vii_dim7}
	La sensibile (Sol diesis) sale alla tonica, la quinta diminuita (Re) e la settima (Fa) scendono per grado.
	\musica{es_cap13_dim_la_1}
	\musica{es_cap13_dim_la_2}
	Nei rivolti il basso segue il suo movimento obbligato: nel \sigla{VII}{6}{5} (basso Si) scende a La, nel \sigla{VII}{4}{3} (basso Re) scende a Do, nel \sigla{VII}{2}{} (basso Fa, la settima) scende a Mi sul \sigla{i}{6}{4}. La sensibile sale sempre al La.
\end{esempio}

I rivolti ricalcano quelli della settima di sensibile (terza specie sul VII), con una sola eccezione, per il \textbf{terzo rivolto}.

\begin{osservazione}[Terzo rivolto senza preparazione]
	\label{oss:dim7_terzo_rivolto}
	A differenza della terza specie, nella quinta specie il \textbf{terzo rivolto} può essere disposto \textbf{senza alcuna necessità di preparare la settima al basso}. Motivo teorico: l'urto di \textbf{seconda aumentata} che si crea tra la fondamentale (sensibile) e la settima è \textbf{enarmonico} all'intervallo consonante di terza minore, e all'udito risulta privo della durezza tipica delle seconde.
\end{osservazione}

Oltre che sul VII grado, questa specie si usa anche come \textbf{accordo alterato sul IV grado} di entrambi i modi (maggiore e minore), come "dominante della dominante" o per scopi coloristici nelle cadenze.

\section{Sesta specie (minore con settima maggiore)}
\label{sec:sesta_specie_dettaglio}

La sesta specie è formata da una \textbf{triade minore} e una \textbf{settima maggiore}. È un accordo \textbf{secondario e artificiale}, di ambito moderno: non viene usato nell'armonia diatonica tradizionale e classica, dove era considerato un accordo alterato.

\begin{regola}[Grado e risoluzione della sesta specie]
	\label{reg:sesta_specie}
	La sesta specie si trova sul \textbf{I grado della scala minore armonica e melodica}, con sigle $m^{maj7}$ e $m\Delta$ (es. La-Do-Mi-\Soldiesis). Al basso ha la fondamentale (I grado) e risolve \textbf{generalmente sul IV grado}.
\end{regola}

La sua settima \textbf{è la sensibile della tonalità}: non può scendere, quindi \textbf{sale} alla tonica (Sol\diesis{} $\rightarrow$ La) e non si raddoppia. È l'unica specie in cui la settima ha moto ascendente. Va preparata come le altre settime secondarie.

\begin{esempio}[Sesta specie in La minore: \sigla{I}{7}{} $\rightarrow$ IV]
	\label{ex:sesta_specie}
	Nello stato fondamentale e nei rivolti la settima Sol\diesis{} sale a La, che è la quinta del IV (Re-Fa-La). Nel terzo rivolto il basso Sol\diesis{} sale a La e il IV è in secondo rivolto.
	\musica{es_cap13_minmagg_la_1}
	\musica{es_cap13_minmagg_la_2}
\end{esempio}

\section{Settima specie (aumentata con settima maggiore)}
\label{sec:settima_specie_dettaglio}

La settima specie è formata da una \textbf{triade aumentata} e una \textbf{settima maggiore}. È anch'essa un accordo \textbf{secondario e artificiale} di ambito moderno/jazzistico, considerato alterato nella prassi classica.

\begin{regola}[Grado e risoluzione della settima specie]
	\label{reg:settima_specie}
	La settima specie si trova sul \textbf{III grado della scala minore armonica e melodica}, con sigle $maj^7(\sharp5)$ e $+\Delta$. Al basso ha la fondamentale (III grado) e risolve \textbf{generalmente sul VI grado}.
\end{regola}

Ha due note con tendenza: la settima Si \textbf{scende} a La, la quinta eccedente Sol\diesis{} \textbf{sale} a La. Nessuna delle due si raddoppia.

\begin{esempio}[Settima specie in La minore: \sigla{III}{7}{} $\rightarrow$ VI]
	\label{ex:settima_specie}
	Nel primo rivolto (basso Mi) il basso sale a Fa; nel secondo (basso Sol\diesis) sale a La e il VI è in primo rivolto; nel terzo (basso Si) scende a La.
	\musica{es_cap13_ecc_la_1}
	\musica{es_cap13_ecc_la_2}
\end{esempio}

Pur appartenendo al catalogo completo delle quadriadi, queste due specie non trovano applicazione nell'armonia diatonica tradizionale del basso dato.

\section{Quadro riassuntivo}
\label{sec:quadro_specie}

\begin{center}
\small
\begin{tabular}{|l|l|l|l|}
\hline
Specie & Esempio & Risoluzione tipica & Moto della settima \\
\hline
1\textsuperscript{a} dominante & \sigla{V}{7}{} $\rightarrow$ I & settima $\searrow$, sensibile $\nearrow$ & scende \\
2\textsuperscript{a} minore & \sigla{II}{7}{} $\rightarrow$ V, \sigla{i}{7}{} $\rightarrow$ iv & quinta discendente & scende \\
3\textsuperscript{a} semidiminuita & \sigla{VII}{7}{} $\rightarrow$ I, \sigla{II}{7}{} $\rightarrow$ V & tonica (sensibile $\nearrow$) o dominante & scende \\
4\textsuperscript{a} maggiore & \sigla{I}{7}{} $\rightarrow$ IV, \sigla{VI}{7}{} $\rightarrow$ II & quinta discendente & scende \\
5\textsuperscript{a} diminuita & \sigla{VII}{7}{} $\rightarrow$ i & tonica, sensibile $\nearrow$ & scende \\
6\textsuperscript{a} minore-maggiore & \sigla{i}{7}{} $\rightarrow$ iv & quinta discendente & \textbf{sale} (è la sensibile) \\
7\textsuperscript{a} con 5\textsuperscript{a} ecc. & \sigla{III}{7}{} $\rightarrow$ VI & 5\textsuperscript{a} ecc.\ $\nearrow$, settima $\searrow$ & scende (la 5\textsuperscript{a} ecc.\ sale) \\
\hline
\end{tabular}
\end{center}

Tutti gli esempi si trasportano in qualsiasi tonalità mantenendo la stessa condotta delle voci: cambiano soltanto i nomi delle note e le alterazioni in chiave.
